% Part: lambda-calculus
% Chapter: lambda-definability
% Section: lambda-definable-recursive

\documentclass[../../../include/open-logic-section]{subfiles}

\begin{document}

\olfileid{lam}{dfl}{ldr}
\olsection{\usetoken{S}{lambda definable} অপেক্ষকগুলি পুনরাবৃত্ত}

সব আংশিক পুনরাবৃত্ত অপেক্ষক যে শুধু !!{lambda definable} তা-ই নয়;
বিপরীত কথাটিও সত্য। অর্থাৎ সব !!{lambda definable} অপেক্ষকই আংশিক
পুনরাবৃত্ত।

\begin{thm}
  \ollabel{thm:lambda-computable} কোনো আংশিক অপেক্ষক~$f$ যদি
  !!{lambda definable} হয়, তবে সেটি আংশিক পুনরাবৃত্ত।
\end{thm}

\begin{proof}
  আমরা শুধু প্রমাণটির রূপরেখা দিই। প্রথমে $\lambd$-পদগুলির পাটিগণিতায়ন
  করি; অর্থাৎ অনুক্রমের প্রচলিত মৌলিক-সংখ্যার-ঘাত সংকেতায়ন ব্যবহার করে
  $\lambd$-পদগুলিকে নিয়মমাফিক গ্যোডেল সংখ্যা দিই। এরপর একটি আংশিক
  পুনরাবৃত্ত অপেক্ষক $\fn{normalize}(t)$ সংজ্ঞায়িত করি, যা কোনো ল্যাম্বডা
  পদের গ্যোডেল সংখ্যা~$t$-কে আর্গুমেন্ট হিসেবে গ্রহণ করে; স্বাভাবিক রূপ
  থাকলে তার গ্যোডেল সংখ্যা ফেরত দেয়, আর না থাকলে অসংজ্ঞায়িত থাকে। তারপর
  $\fn{toChurch}$ ও~$\fn{fromChurch}$ নামে দুটি আংশিক পুনরাবৃত্ত অপেক্ষক
  সংজ্ঞায়িত করি, যা স্বাভাবিক সংখ্যা থেকে সংশ্লিষ্ট চার্চ সংখ্যাপদের
  গ্যোডেল সংখ্যায় এবং সেই গ্যোডেল সংখ্যা থেকে স্বাভাবিক সংখ্যায় চিত্রিত
  করে।

  এই পুনরাবৃত্ত অপেক্ষকগুলি দিয়ে~$f$-কে একটি আংশিক পুনরাবৃত্ত অপেক্ষক
  হিসেবে সংজ্ঞায়িত করা যায়। কোনো $\lambd$-পদ~$F$, $f$-কে
  !!{lambda define}s। $f(n_1, \dots, n_k)$ গণনা করতে প্রথমে
  $\fn{toChurch}(n_i)$ ব্যবহার করে সংশ্লিষ্ট চার্চ সংখ্যাপদগুলির গ্যোডেল
  সংখ্যা বার করো; এগুলি $\Gn{F}$-এর সঙ্গে জুড়ে
  $F \num{n_1}\dots\num{n_k}$ পদের গ্যোডেল সংখ্যা পাও। এবার ওই গ্যোডেল
  সংখ্যার উপর $\fn{normalize}$ ব্যবহার করো। $f(n_1, \dots, n_k)$
  সংজ্ঞায়িত হলে $F \num{n_1}\dots\num{n_k}$-এর একটি স্বাভাবিক রূপ
  থাকে (যেটি অবশ্যই একটি চার্চ সংখ্যাপদ), আর অন্যথায় তার কোনো স্বাভাবিক
  রূপ থাকে না (তাই
  \[\fn{normalize}(\Gn{F\num{n_1}\dots\num{n_k}})\]
  অসংজ্ঞায়িত)। সবশেষে স্বাভাবিকীকৃত পদটির গ্যোডেল সংখ্যার উপর
  $\fn{fromChurch}$ ব্যবহার করো।
\end{proof}

\end{document}
